\ifdefined\isincluded
\else
\documentclass[a4paper, 12pt]{report}
\usepackage[margin=0.8in]{geometry}
\usepackage{amsmath}
\usepackage{amsfonts}
\usepackage{amssymb}
\usepackage{graphicx}
\usepackage[english]{babel}
\usepackage[autostyle, english = american]{csquotes}
\MakeOuterQuote{"}
\usepackage{setspace}
\onehalfspacing
\usepackage{booktabs}
\usepackage{longtable}
\usepackage{tabularx}
\usepackage{array}
\usepackage{url}
\usepackage{cite}

\makeatletter
\let\oldincludegraphics\includegraphics
\renewcommand{\includegraphics}[2][]{%
  \IfFileExists{#2}{%
    \oldincludegraphics[#1]{#2}%
  }{%
    \IfFileExists{images/#2}{%
      \oldincludegraphics[#1]{images/#2}%
    }{%
      \fbox{\parbox[c][3cm][c]{0.8\linewidth}{\centering \texttt{[Image: #2]}}}%
    }%
  }%
}
\makeatother

\begin{document}
\chapter{Taxonomy of Enterprise Agent Failure Modes and Diagnostics}
\fi

\label{chap.7}
\section{Taxonomy of Enterprise Agent Failure Modes}
To understand why agents fail and how failures relate to underlying latent geometry, we examined all 114 benchmark tasks and categorized results into four mutually distinct operational states:
\begin{enumerate}
    \item \textbf{Clean Success (No Error):} Ground truth is reflected in all final customer outputs, policy validations, and database changes.
    \item \textbf{Tool Selection Error:} The agent called tools out of topological order, passed corrupted or hallucinated JSON parameters, or selected an incorrect tool name.
    \item \textbf{Premature Termination:} The agent finished execution before verifying all crucial limitations or completing all required changes.
    \item \textbf{Policy Constraint Violation:} The agent changed shipped orders or refunded non-refundable items, two actions that were specifically prohibited by business policy.
\end{enumerate}

Table~\ref{tab:error_profiles} shows the empirical failure distributions' geometric profiles.

\begin{table}[h!]
\centering
\small
\caption{Geometric Profiles for Text and Latent ReAct Across Failure Modes.}
\label{tab:error_profiles}
\vspace{3pt}
\begin{tabular}{lcccccc}
\toprule
\textbf{Architecture \& Mode} & \textbf{Count} & \textbf{$D_1$} & \textbf{$D_2$} & \textbf{$D_3$} & \textbf{Decay Slope} & \textbf{$D_{ij}$} \\
\midrule
\textbf{Latent ReAct} & & & & & & \\
Success (No Error) & 33 & 0.0672 & 0.0196 & 0.0199 & -0.0237 & 51.40 \\
Tool Selection Error & 54 & 0.0626 & 0.0186 & 0.0194 & -0.0216 & \textbf{51.60} \\
Premature Termination & 17 & 0.0698 & 0.0162 & 0.0190 & -0.0254 & 51.39 \\
Policy Violation & 10 & 0.0676 & 0.0149 & \textbf{0.0164} & \textbf{-0.0256} & 51.49 \\
\midrule
\textbf{Text ReAct} & & & & & & \\
Success (No Error) & 30 & 0.0388 & 0.0331 & 0.0362 & -0.0013 & 10.60 \\
Tool Selection Error & 53 & 0.0365 & 0.0310 & 0.0334 & -0.0016 & 10.69 \\
Premature Termination & 18 & 0.0423 & 0.0395 & 0.0330 & -0.0046 & 10.49 \\
Policy Violation & 13 & 0.0352 & 0.0261 & 0.0279 & -0.0037 & 10.71 \\
\bottomrule
\end{tabular}
\end{table}

\section{Statistical Correlation: Latent Geometry vs. Error Triggers}
We computed Pearson correlation coefficients between task-level geometric measures and the frequency of errors. Key findings include:

\begin{itemize}
    \item \textbf{Planner-Tool Divergence Triggers Tool Errors ($r = +0.212, p = 0.0238$):} 
    The inter-agent distance $D_{12}$ between the Planner ($A_1$) and Tool-Caller ($A_2$) has a positive and substantial correlation with Tool Selection Errors. When $D_{12}$ is larger than $51.6$, the Tool-Caller drifts away from the Planner's semantic subspace, resulting in tool hallucinations.
    \item \textbf{Policy Violations are Confident False Attractors ($r = -0.198, p = 0.0351$):} 
    Policy violations exhibit the lowest final drift ($D_3 = 0.0164$) and the highest decay slope ($-0.0256$). This shows that policy violations are not caused by miscommunication or hesitation, but rather by the agent rapidly converging into a false attractor basin (e.g., confidently misinterpreting return rules).
    \item \textbf{Information Deficit Causes Premature Termination:} 
    Premature termination is associated with low initial effective rank ($\text{ER}_{r=1} < 89, r = -0.185, p = 0.0489$), indicating that the model was not able to fully capture relational connections during round $r=1$.
\end{itemize}

Figure~\ref{fig:error_modes} shows the geometric characteristics for every error mode.

\begin{figure}[h!]
\centering
\includegraphics[width=0.9\linewidth]{error_modes_vs_geometry.png}
\caption{Subspace separation, inter-agent distance, and representation drift distribution among different error modes are examples of geometric failure diagnostics.}
\label{fig:error_modes}
\end{figure}
\section{Early Predictive Latent Probing}
To determine if task failure might be detected prior to completion of execution, we trained a linear probe on the returned latent representation at Round 1 ($H_{r=1}$).
\begin{itemize}
    \item \textbf{Probe Architecture:} A single linear projection layer ($\mathbb{R}^{2560} \to \mathbb{R}^2$) was trained using binary cross-entropy.
    \item \textbf{Result:} The probe predicts a trajectory's success or failure at Round 3 with \textbf{71.4\% accuracy} based solely on its $r=1$ latent state.
\end{itemize}
By verifying that failure signatures show up in representation geometry early in consideration, this diagnostic feature creates new opportunities for adaptive runtime rerouting.

\section{Resolving the Tool Selection Trade-Off}
Latent ReAct had an accuracy of 54.2\% in raw isolated tool selection, compared to 57.5\% for Text ReAct. This occurs as a result of Text ReAct's explicit tool names in its prompt headers, which provide a superficial lexical matching advantage. However, because text deliberation cannot self-correct relational errors, Text ReAct's overall accuracy on Hard tasks drops to 45.9\%. Conversely, Latent ReAct uses recursive latent feedback to recover from early tool errors and achieves 50.8\% final completion accuracy.

\ifdefined\isincluded
\else
\end{document}
\fi
